% Build the layer, Chapter 8, first build: a witness log.
% Inserted in \S\ref{sec:fh-authority}, after \S\ref{sec:fh-witness}, which
% states what a witness log must do and ends with an example. This section
% builds the part of the log that the example needs, in the kernel's
% decision core, and runs the example.
%
% No code or output in this file is typed by hand. Each listing is a file
% under excerpts/, which the kernel's tools/excerpts.mjs cuts from its
% source and its recorded runs (see figures/pd-listing.tex). To change a
% listing, change the kernel and cut again. The prose refers to functions
% by name and never by line number, so a listing that moves within its file
% does not make a sentence wrong.
%
% \pdlistingroom{n} before a listing of n lines keeps it on one page. When a
% listing changes length, change the number with it.
\subsection{Build the layer: a witness log}
\label{sec:fh-log}

\noindent\emph{Express lane: Acme's harbor signs two roots for epoch~8
and shows each to a different auditor. A log that stores the latest root
it is given loses the first one, and then nobody can show what happened. The log
built here keeps the first root and holds a second one beside it. The two
signed roots are the evidence, and anyone can check them. Two logs that
each hold one of the roots get the evidence when they exchange what they
hold. An exhaustive test of three logs reaches 8,000 states. In 4,608 of
them two logs hold different roots, and in every one of those a single
exchange leaves both logs with the evidence.}

Design invariant~\ref{prop:fh-witness} lists five requirements for a
witness log, and the example after it ends with an equivocation: Acme signs
two roots for epoch~8. None of that has been run. This section builds the
part of the log that the example needs, in \pdkernelname{}, and runs the
example against it.

The cast is the example's. Acme's harbor has recorded four card events by
the end of epoch~7. In epoch~8 it records a fifth, which revokes the card
$C_A$. Beta and the staging server each run an auditor, a program that
reads Acme's root at a log and checks what it reads. Wren is an outsider,
as she was in the builds of \pdchapref{anchor}{The Anchor Protocol}. She
runs no harbor and holds no harbor's key.

The section builds two of the five requirements: a log records one root
for each epoch, and two signed roots are evidence of equivocation. It does
not build the proofs that tie one epoch's root to the next. The boundary at
the end lists what else is missing.

\subsubsection{A root, and the text a harbor signs}

A harbor's root is one hash over its whole list of events. The decision
core computes it (\weld{prereq:sec:prereq-build}{decision core}), so a
harbor and an auditor get the same value from the same list
(Listing~\ref{alg:fh-log-root}).

\pdlistingroom{32}
\pdexcerpt[caption={The root of a list of events.},label={alg:fh-log-root}]{ch08-root}

\texttt{tree} is the Merkle tree hash of Certificate
Transparency~\pdcite{laurie2013ct}. It splits the list in two and hashes
each part, and the hash of those two results is the root. For Acme's first
four events that gives the tree of the example,
$H\bigl(H(h_1,h_2),H(h_3,h_4)\bigr)$. With five events the split falls
after the fourth, so the root for epoch~8 is $H(R_A^{(7)},h_5)$, as the
example has it. The code adds one byte to the example's formula: a leaf is
hashed after a 0 and a pair after a 1. Without that byte, a harbor could
record the two hashes of a pair as if they were one event, and two
different lists would have the same root. A test in the kernel compares
\texttt{tree} with the known roots of Certificate Transparency's
reference test inputs.

The root of Acme's four events begins \texttt{38936ba2}, and the root of
all five begins \texttt{3c657ada}. The runs below print those eight
digits.

A harbor signs more than the root (Listing~\ref{alg:fh-log-signed}).

\pdlistingroom{17}
\pdexcerpt[caption={The text a harbor signs, and the check of its
signature.},label={alg:fh-log-signed}]{ch08-signed}

The signed text holds the harbor's public key and the epoch as well as the
root. A log keeps no table of harbors. A harbor's key is its name, and
\texttt{signed} verifies the signature with that key. The epoch is in the
text for a reason that the verification below returns to.

\subsubsection{Naive: keep the latest root}

A log is a table, and the obvious table stores each signed root under its
harbor and epoch (Listing~\ref{alg:fh-log-latest}). Its \texttt{publish}
checks the signature before it stores anything, so only Acme can store a
root as Acme's. With a harbor that signs one root for each epoch, this log
works. For epoch~7 an auditor reads the root \texttt{38936ba2} and asks
Acme for its events. If the four events hash to that root, the auditor
knows what Acme had recorded by then.

\pdlistingroom{12}
\pdexcerpt[caption={The first attempt. A root is stored under its harbor
and epoch.},label={alg:fh-log-latest}]{ch08-latest}

\subsubsection{Break: a second root for the same epoch}

Acme publishes its root for epoch~8, and Beta's auditor reads it. Then
Acme signs a second root for epoch~8, over its first four events only, and
publishes that one too (Listing~\ref{alg:fh-log-run-latest}).

\pdlistingroom{22}
\pdrun[caption={A second root for one epoch, at the first log. The word at
the left of each step says who acts.},label={alg:fh-log-run-latest}]{two-roots-latest}

\texttt{publish} stored the second root where the first had been. Both
auditors did their job. Each checked a signature that verifies, and each
was shown events that hash to the root it read. Beta's auditor concluded
that $C_A$ is revoked, and the staging server's that it is not. So the
staging server goes on honoring a card that Acme has revoked.

The log no longer holds the first root. Beta's auditor read it and kept no
copy, so at step~5 it is given the second root like everyone else, and
nobody can now show that Acme ever signed the first. This is the
equivocation that \S\ref{sec:fh-scope} describes: one party signs two
versions of a record and shows each to a different audience.

\subsubsection{Fix: keep the first root, and hold the second as evidence}

The second requirement of Design invariant~\ref{prop:fh-witness} is that a
log records at most one root for a harbor and an epoch. In
\pdkernelname{} the log hands its contents to the decision core with each
root, and the core decides (Listing~\ref{alg:fh-log-publish}).

\pdlistingroom{21}
\pdexcerpt[caption={A root is published to a
log.},label={alg:fh-log-publish}]{ch08-publish}

The first root that arrives for a harbor and an epoch is recorded. The
same root arriving again is \texttt{Known} and changes nothing, so a
harbor that is unsure its root arrived can send it again. A different
root goes to \texttt{second} (Listing~\ref{alg:fh-log-second}).

\pdlistingroom{20}
\pdexcerpt[caption={A second root for an epoch. A line marked
\texttt{MUTANT} is compiled only into a test build of the core that has
that guard removed. The verification below runs each of those
builds.},label={alg:fh-log-second}]{ch08-second}

\texttt{second} leaves the recorded root where it is and stores the two
roots as a pair. The requirement's words are that a log ``refuses a
second'' root. The core does refuse to make it the root for that epoch,
and it still keeps the second root, because the pair is the proof
(Listing~\ref{alg:fh-log-run-core}). Exercise~\ref{ex:fh-log-third} asks
what a third root does.

\pdlistingroom{14}
\pdrunlines[caption={The same run at the kernel's log. The steps left out
are the first two of
Listing~\ref{alg:fh-log-run-latest}.},label={alg:fh-log-run-core}]{two-roots-core}{1,10-22}{1 and 10\textendash 22}

Both auditors now get the same answer, and the answer is the pair. The
fourth requirement is that an auditor who holds two differently signed
roots for one harbor and epoch can verify the contradiction at once.
\texttt{proves} is that check (Listing~\ref{alg:fh-log-evidence}).

\pdlistingroom{21}
\pdexcerpt[caption={The check of evidence. It reads the two signed roots
and no log.},label={alg:fh-log-evidence}]{ch08-evidence}

An auditor does not have to trust the log that showed it the pair.
\texttt{proves} is given no log. It makes four checks: both signatures
verify, both roots name one harbor, both are for one epoch, and the two
roots differ. It returns the harbor and the epoch, and the auditor
compares those with the harbor and epoch it asked the log about.

\subsubsection{Break again: two logs, and a root for each}

One log is a single point of failure, and \S\ref{sec:fh-not} asks for a
few logs that anyone can read. So Acme changes its approach. It publishes
one root to each of two logs (Listing~\ref{alg:fh-log-run-separate}).

In that run each log does what the fix requires. Each holds one root for
epoch~8 and would keep a second as evidence, but neither is given a
second. The auditors end up split exactly as they were at the first log,
and this time both logs obeyed the rule.

Design invariant~\ref{prop:fh-witness} expects this. Several logs cannot
promise to hold the same root unless they agree among themselves, and
agreeing is consensus, which this design does not use. A log cannot refuse
a root because of what some other log holds.

\pdlistingroom{20}
\pdrunlines[caption={One root to each of two logs. Both logs ask the core.
Steps 6 and 7, left out, repeat steps 3 and
4.},label={alg:fh-log-run-separate}]{two-logs-separate}{1-17,24-26}{1\textendash 17 and 24\textendash 26}

\subsubsection{Fix: the logs exchange the roots they hold}

What the logs can do is compare. Each sends the other every signed root it
holds, and each publishes what it receives to itself
(Listing~\ref{alg:fh-log-exchange}).

\pdlistingroom{8}
\pdexcerpt[caption={An exchange between two
logs.},label={alg:fh-log-exchange}]{ch08-exchange}

Nothing new has to be trusted. A root that arrives from another log goes
through \texttt{publish} like a root that arrives from a harbor, and it is
recorded only if its harbor signed it. A log that wanted to accuse Acme
falsely would have to forge Acme's signature. Gates in
\pdchapref{anchor}{The Anchor Protocol} treat a notice of a revoked card
the same way (\weld{anchor:sec:anchor-build-told}{notices}): the receiver
checks it as if it had been presented there.

A log sends the second root of each pair it holds, as well as the roots it
has recorded. A third log that has seen neither of Acme's roots then gets
the evidence from either of the first two.

\pdlistingroom{12}
\pdrunlines[caption={The same run with logs that exchange. The steps left
out are the first four of
Listing~\ref{alg:fh-log-run-separate}.},label={alg:fh-log-run-exchange}]{two-logs-core}{1,16-26}{1 and 16\textendash 26}

Log~1 still holds the root it recorded first, and so does log~2. Each now
holds the other's root as the second of a pair, and each auditor checks
the pair for itself.

The exchange happened at step~5. The staging server's auditor read that
$C_A$ is not revoked at step~4. The log finds the equivocation after the
fact and prevents nothing, and \S\ref{sec:fh-attacks-in} is plain about
that: a bond on a witness pays out once an equivocation is confirmed, and
does not make the proof arrive by any deadline.

\keyidea[Keep the first root, and keep the second beside it]{A log cannot
stop a harbor from signing two roots for one epoch. It can refuse to
replace the first, and hold the second as evidence. The harbor's two
signatures prove the equivocation to anyone, and no witness has to be
believed.}

\Needspace{12\baselineskip}
\subsubsection{Verify: an outsider, then every order}

\paragraph{An outsider.}
Evidence that is easy to check must also be hard to fake. In
Listing~\ref{alg:fh-log-run-outsider} Acme signs one root for each epoch,
and Wren tries to make it look as if Acme signed two.

\pdlistingroom{36}
\pdrun[caption={An outsider at the kernel's log. Beta's harbor publishes
a root of its own at step
3.},label={alg:fh-log-run-outsider}]{false-evidence-core}

Steps 6 to 9 are the four checks of \texttt{proves}, each refusing a pair
that fails it. Step~5 is the retry: Acme's own root, published again, is
not a second root. Step~4 is refused by \texttt{signed}.

Step~7 is worth a second look. Acme's root for epoch~7 covers four events.
The second root that Acme signed for epoch~8, in the runs above, covers
the same four events, so it is the same hash. If a harbor signed its root
and nothing else, Acme's honest signature for epoch~7 would also be a
valid signature on a second root for epoch~8, and Wren could publish it
as one. That is why the epoch is inside the signed text of
Listing~\ref{alg:fh-log-signed}. \emph{Now you try:}
Exercise~\ref{ex:fh-log-epoch} works the forgery through.

\paragraph{Every pair.}
The kernel offers the core every pair that can be made from five signed
roots (Listing~\ref{alg:fh-log-pairs}). Acme signed three of them: one
for epoch~7, and two different ones for epoch~8. Beta signed one for its
own epoch~8. Wren signed the fifth, which names Acme.

\pdlistingroom{20}
\pdrun[caption={Every pair of five signed roots, offered as
evidence.},label={alg:fh-log-pairs}]{evidence}

Of the 25 pairs the core accepts 2, which are the two orders of Acme's
two roots for epoch~8. The pair of \texttt{acme-7} and \texttt{acme-8'}
is refused as \texttt{epoch}, although their roots are the same hash.

\paragraph{Every order.}
Which root a log records depends on which arrives first, so a few runs
are not enough. The exhaustive test takes three logs and the same five
roots. Any root may be published to any log, as often as you like and in
any order. Any two logs may also exchange
(Listing~\ref{alg:fh-log-walk}). The test stops when no event leads to a
state it has not seen.

In every state the test compares each log with a small reference model.
The model checks no signature. It is told which roots are real, and it
keeps its own account of what has reached each log.

\pdlistingroom{30}
\pdrun[caption={Every sequence of events for three logs, at each
design.},label={alg:fh-log-walk}]{logs}

The first log drops a root in 20,224 of its 21,952 states. The logs that
ask the core never drop one. They also never hold a root that its harbor
did not sign, or evidence against a harbor that signed one root for the
epoch. Those three columns are zero for both.

The column \texttt{split} is not zero for any design, and it cannot be.
In 4,608 of the core's 8,000 states two logs hold different roots for one
epoch, because Acme can always give two logs two roots. The last column
is what the exchange adds. Without it, one of the two logs lacks the
evidence in 3,840 of those states. With it, one exchange between the two
logs leaves both holding the evidence, in all 4,608. \emph{Now you try:}
Exercise~\ref{ex:fh-log-count} counts the same columns by hand for two
logs and two roots.

\paragraph{Each guard removed.}
The build has six guards. The core can be compiled with any one of them
removed, and the checks should then fail
(Listing~\ref{alg:fh-log-guards}).

\pdlistingroom{31}
\pdrun[caption={The exhaustive test for two logs, on the core and on six
builds of it with one guard removed from
each.},label={alg:fh-log-guards}]{logs-guards}

Every row differs from the first. A core that records a root whatever its
signature holds Wren's root in 2,736 states. A core that lets a later root
replace the first is the first log again. The three builds that weaken
\texttt{proves} alone leave the logs untouched, and each accepts pairs
that are not evidence. The last row is the retry: a core that takes one
root, sent twice, for two roots holds evidence against a harbor that did
nothing wrong in 5,777 of its 6,561 states. This table uses two logs
because a core with a guard removed reaches far more states than the
real one.

\paragraph{The check, as a protocol.}
The exhaustive test uses five roots that the kernel made. ProVerif asks
the same question about any roots an attacker can make
(\weld{anchor:sec:correct}{ProVerif}). In the model, Beta signs one root
for each epoch, and Acme signs whatever it is asked to. The attacker sees
every signed root and has a key of its own
(Listing~\ref{alg:fh-log-models}).

\pdlistingroom{25}
\pdrun[caption={ProVerif~2.05 on the check of evidence, and on five
controls.},label={alg:fh-log-models}]{models-witness}

\texttt{evidence.pv} is the check as the core makes it. Its first result
says that each root an auditor accepts was signed by the harbor, for that
epoch. Its second says that no auditor ever accepts evidence against
Beta. Its third is \texttt{false} on purpose: an auditor can accept
evidence against Acme, so the first two results are not true merely
because nothing is ever accepted.

The other five files are controls. Four leave out one check of
\texttt{proves} each, and the fifth takes the epoch out of the signed
text. In every one of them an auditor accepts evidence against Beta.

\pdscope{The exhaustive test covers three logs and the five roots printed
in Listing~\ref{alg:fh-log-pairs}. It counts states and says nothing
about how likely a state is, or how long an exchange takes to happen. The
ProVerif models cover the check of two signed roots, and not the log that
stores them.}

\begin{pdboundary}
\setlength{\parskip}{4pt plus 1pt}
Nothing here ties one epoch to the next. Acme can publish a root for
epoch~9 over a list that leaves out the revocation of $C_A$, and the log
records it, because it is the first root for epoch~9. The third
requirement of Design invariant~\ref{prop:fh-witness} answers this with a
consistency proof, and the kernel does not build that proof yet
(Exercise~\ref{ex:fh-log-rewrite}).

An auditor trusts the log to show the evidence it holds. A log that keeps
a pair back shows one signed root, and the auditor cannot tell. The
design invariant has witnesses sign the heads of their own logs, so that
a log which shows two auditors different contents can be caught in turn.
A log here signs nothing (Exercise~\ref{ex:fh-log-withhold}).

The exchange has no schedule. Two logs that never exchange never get the
evidence, and \S\ref{sec:fh-lim-equiv} explains why the chapter gives no
bound on how long the evidence takes to arrive.

A harbor is its public key, and any key can publish. The joining ceremony
of \S\ref{sec:fh-admission} decides which keys belong to members, and it
is not built. Nothing is paid out on evidence either. The bond that a
harbor or a witness would lose belongs to the settlement of
\S\ref{sec:fh-settlement}.

The logs run in one process. Running each harbor and each log as its own
process is a later build, and so is the transfer ceremony of
\S\ref{sec:fh-xfer}.
\end{pdboundary}
